\begin{textsample}{NEG Dim 6 – vem\_ed\_06 – Score 1.00 – vem\_ed\_06/t022.txt}  \label{ex:f6_neg_084}
QUEM É QUEM Lavern Samuels é presidente do Directors Forum of the International Education Association of South Africa (Fórum de Diretores da IEASA) e diretor do escritório internacional da Durban University of Technology (DUT, África do Sul).
A DUT tem cerca de 33 mil alunos em seis faculdades e recebeu em fevereiro de 2021 o Global Award for Innovation and Excellence in Internationalisation (Prêmio Global pela Inovação e Excelência na Internacionalização), concedido pela Association of International Education Administrators (AIEA), que fica nos EUA.
Foi a primeira vez que uma universidade fora da Europa e dos Estados Unidos ganhou essa honraria.
A seguir, uma breve entrevista com Samuels, concedida a VEm por e-mail.
Conte-nos um pouco sobre sua experiência no Fórum de Diretores da IEASA.
Sou o atual presidente do Fórum de Diretores da IEASA.
É um grupo de diretores de todos os escritórios internacionais na África do Sul.
Trabalhamos juntos para compartilhar \textbf{boas} práticas e expertise.
Nós nos apoiamos no trabalho da internacionalização da educação e ajudamos a construir capacitações.
Qual a importância das colaborações Sul-Sul?
As colaborações Sul-Sul são importantes para fortalecer o sistema de educação superior no Hemisfério Sul.
Trata-se de trazer as vozes do Sul Global para a mesa, promover inclusão e diversidade e fornecer uma plataforma para que suas narrativas únicas sejam compartilhadas, apreciadas e valorizadas.
Também se trata de reforçar o que nos une e moldar a agenda internacional.
Atualmente, quantos projetos COIL (Collaborative Online International Learning) estão ativos na DUT?
Quantos alunos e professores participam?
Temos cerca de 30 projetos ativos, aproximadamente 50 professores e 1000 alunos participando nos projetos.
Quais são os planos para o projeto COIL@DUT, lançado em novembro de 2020?
Queremos que cada curso tenha um projeto COIL e que todos os estudantes da universidade participem pelo menos uma vez durante sua formação.
PCIs em debate Em 31 de março, uma mesa-redonda sobre os Projetos Colaborativos Internacionais (PCIs) encerrou o Seminário Internacional de Tecnologia, Educação e Sociedade, da Fatec Itaquaquecetuba.
Após a abertura, por Neusa Haruka Sesaki Gritti, palestrou Osvaldo Succi Jr., coordenador dos PCIs.
Houve ainda um relato de experiência da professora Paula Pudo e uma breve apresentação de Patrícia Patrício, sobre a comunicação dos PCIs.
A mediação foi de Fábio Barbosa Lima, da área de políticas linguísticas da ARInter.

% matched lemmas: bom
\end{textsample}
